\chapter{Literature Review}
\label{Chapter2}

\section{Secure Boot and Boot-Chain Protection}
Secure boot anchors trust in immutable first-stage code (BootROM/ROM-resident loader) that
authenticates subsequent software before execution. Early work on measured boot established
trust properties over commodity boot stacks~\cite{parno2010bootstrapping}, while later surveys
summarized secure-boot techniques and practical constraints in embedded
deployments~\cite{rv2018securebootreview}. BootROM designs emphasize validation of system images
at startup~\cite{bin2020bootrom}. As embedded systems incorporate more complex boot sequences,
boot-chain protection extends integrity across multiple software stages and configuration
artifacts~\cite{devic2011fpga}. Device identity composition (DICE) further binds boot trust to
device identity and measurements~\cite{jaeger2020diceharder}, with recent surveys covering
trade-offs across deployment models~\cite{wang2022securebootsurvey}. These lines of work provide
boot-time authenticity but leave post-boot debug and programming access less constrained.

\section{Debug Interface Security}
Debug and test access are high-impact attack surfaces because they provide invasive control over
execution, memory, and device state. Early work extended JTAG for at-speed debug and emulator
implementations~\cite{vandelogt2003jtag,chen2008jtagemu}, while IEEE~1149.1 standardized the test
access port and boundary-scan architecture~\cite{ieee1149_1}. Boundary-scan/test modes can be
abused when an adversary enters privileged states~\cite{ali2014boundaryscanattack}. More recent
work adds explicit authentication to debug access, including JTAG authentication IP and
identity-based schemes~\cite{lapeyre2022jtagauth,wang2020schnorrjtag,walimbe2014jtagaxi}. These
approaches harden the interface but are often external to the execution model.

\section{Trusted Execution and Minimal TCB Designs}
Hardware-enforced trusted execution evolved from secure-processor and minimal-TCB designs such as
AEGIS and minimal TCB code execution~\cite{suh2005aegis,mccune2007minimaltcb}. Measured-execution
systems such as Flicker, TrustVisor, and Memoir reduced trusted code while preserving continuity
guarantees~\cite{mccune2008flicker,mccune2010trustvisor,parno2011memoir}. TrustLite targeted
constrained devices~\cite{koeberl2014trustlite}, and TrustZone-style partitioning became widely
deployed, with evaluations highlighting attack surfaces in heterogeneous
SoCs~\cite{zhang2017trustzoneeval}. Open-platform efforts such as Keystone demonstrate TEEs with
explicit, auditable interfaces~\cite{lee2020keystone}. Surveys provide broader context on enclave
designs and assumptions~\cite{costan2017secureprocessors,ahmed2025riscvtee}, while software-root-of-
trust work explores small, well-defined bases of trust~\cite{gligor2019establishing}.

\section{Runtime Integrity and Attestation}
Beyond boot-time authenticity, runtime integrity mechanisms provide evidence of correct execution.
Cumulative attestation kernels extend attestation by accumulating evidence across
execution~\cite{lemay2012cumulative}, followed by remote attestation for embedded
systems~\cite{kylanpaa2016remoteattest} and collective attestation across embedded
networks~\cite{ibrahim2018collectiveatt}. Control-flow attestation progressed from software-based
approaches (C-FLAT) to hardware-assisted designs (LO-FAT) and later log-based and
micro-architectural variants~\cite{abera2016cflat,dessouky2017lofat,liu2019logbasedcfa,benyehuda2022nanovised}.
Formal verification provides high-assurance evidence for critical components, as demonstrated by
seL4~\cite{klein2009sel4,klein2014tocs}.

\section{Software-Only Hardening}
Software protection techniques deter tampering and reverse engineering by transforming code or
embedding secret-dependent checks. Control-Flow Integrity constrains indirect control
transfers~\cite{abadi2005cfi}, while return-oriented programming demonstrates code-reuse attacks
that motivate defenses for sensitive execution paths~\cite{checkoway2010rop}. Recent taxonomies
summarize layered obfuscation techniques and security goals~\cite{tran2020layered}. These methods
provide valuable context but do not directly enforce device-bound authorization.

\section{Toolchain Infrastructure and ISA Extensibility}
Reliable low-level enforcement requires toolchain support. LLVM provides a foundation for program
analysis and transformation~\cite{lattner2004llvm}. Emscripten and LLBT demonstrate LLVM-based
retargeting and binary translation~\cite{zakai2011emscripten,shen2012llbt}. RISC-V’s extensibility
supports custom instruction integration with conventional toolchains~\cite{waterman2014riscv}.
These efforts motivate ISA-level authorization mechanisms that are emitted and preserved by the
compiler and binary utilities.

\section{Summary of Research Gaps}
Existing approaches provide strong foundations for boot integrity, debug hardening, trusted
execution, and runtime attestation. However, many mechanisms remain external to the execution
model or rely on host-side policy. The literature indicates a need for device-local authorization
that is invoked as part of execution, is bound to compiler-emitted artifacts, and keeps key
material outside the normal memory hierarchy. This gap motivates the proposed ISA-triggered,
key-gated authorization mechanism.
